\rtmtight
\begin{tblr}{width=\textwidth,colspec={|Q[c,m,2.1cm]|X[l,m]|},hlines,vlines,hspan=minimal,rowsep=1.5pt,colsep=3pt,row{1}={bg=headgray,font=\bfseries}}
\SetCell{c,m}\hfil\logo\hfil & \textbf{POLITEKNIK NEGERI MALANG}\par\textbf{JURUSAN TEKNOLOGI INFORMASI}\par\textbf{PROGRAM STUDI: D4 TEKNIK INFORMATIKA} \\
\SetCell[c=2]{l} \textbf{RENCANA TUGAS MAHASISWA (RTM) DAN RUBRIK PENILAIAN} & \\
Nama Mata Kuliah & Konsep Teknologi Informasi \\
Kode Mata Kuliah & RTI251002 \quad \textbf{sks / Jam perminggu:} 2 SKS/4 jam \quad \textbf{Semester:} 1 \\
Dosen Pengampu & \begin{enumerate}\item Ariadi Retno Ririd, S.Kom., M.Kom.\item Budi Harijanto, S.T., M.MKom.\item Deddy Kusbianto PA, Ir., M.MKom.\item Banni Satria Andoko, S.Kom., M.MSI., Dr.Eng.\end{enumerate} \\
\end{tblr}
\assessmentblock{Tugas Rangkuman dan Studi Kasus TI}{Tugas individu dan kelompok}{SCPMK0901-00201--SCPMK0208-00204}{Mahasiswa mengerjakan tugas mingguan berupa rangkuman materi, ulasan topik, dan latihan soal sesuai bahan kajian tiap pertemuan.}{Merangkum materi atau mengulas topik yang ditugaskan, mengerjakan latihan soal (representasi data, aljabar Boolean, flowchart), mendiskusikan hasil dalam kelompok, dan mengumpulkan melalui LMS.}{Dokumen rangkuman/ulasan, jawaban latihan soal, dan bahan diskusi kelompok.}{Ketepatan konsep dan isi rangkuman/ulasan & 10 & Rangkuman menjelaskan konsep, fungsi, dan contoh secara akurat sesuai materi pekan berjalan dengan sumber yang jelas. & Konsep utama benar, tetapi sebagian penjelasan dangkal atau sumber rujukan tidak lengkap. & Rangkuman keliru, menyalin tanpa pemahaman, atau tidak sesuai topik yang ditugaskan. \\ \\
Ketepatan penyelesaian latihan soal dan studi kasus & 9 & Soal representasi data, aljabar Boolean, flowchart, dan ulasan kasus diselesaikan runtut dengan langkah dan hasil yang benar. & Sebagian besar langkah benar, tetapi terdapat kesalahan perhitungan atau simbol yang minor. & Jawaban salah arah, tidak menunjukkan proses pengerjaan, atau tidak selesai. \\ \\
Keaktifan diskusi dan kejelasan penyajian tugas & 7 & Aktif bertanya dan menjawab dalam diskusi kelompok; tugas disajikan jelas, lugas, dan dikumpulkan tepat waktu. & Cukup aktif berdiskusi, tetapi penyajian tugas kurang rapi atau pengumpulan terlambat. & Pasif dalam diskusi dan penyajian tugas sulit dipahami atau tidak dikumpulkan. \\}

\assessmentblock{Kuis 1 Konsep Dasar TI}{Kuis (Ujian Online)}{SCPMK0901-00201}{Mahasiswa mengerjakan kuis online tentang konsep teknologi, inovasi teknologi, perkembangan IPTEK, dan etika rekayasa.}{Mengerjakan soal kuis secara mandiri melalui platform ujian online dalam waktu 60 menit.}{Jawaban kuis online yang terekam pada platform ujian.}{Penguasaan konsep teknologi informasi dan inovasi teknologi & 5 & Menjawab benar soal konsep TI, fungsi, pengelompokan, komponen TI, serta membedakan inovasi SI dan TI modern dengan contoh tepat. & Sebagian jawaban benar, tetapi masih tertukar antara konsep TI dan contoh inovasinya. & Jawaban keliru atau hanya menebak tanpa menunjukkan pemahaman konsep. \\ \\
Pemahaman perkembangan IPTEK dan dampaknya & 4 & Menjelaskan perkembangan IPTEK, contohnya di bidang pendidikan, serta dampak dan solusinya dengan tepat. & Menyebutkan perkembangan IPTEK dengan benar, tetapi dampak atau solusinya kurang tepat. & Tidak mampu mengaitkan perkembangan IPTEK dengan dampak dan solusinya. \\ \\
Pemahaman etika rekayasa dan etika profesi TI & 3 & Mengidentifikasi jenis isu etika TI dan peran etika profesi TI pada kasus yang diberikan dengan tepat. & Mengenali isu etika, tetapi penjelasan peran etika profesi masih umum. & Tidak mampu mengenali isu etika TI pada kasus yang diberikan. \\}

\assessmentblock{UTS Konsep TI dan Sistem Komputer}{UTS (Ujian Online)}{SCPMK0901-00201, SCPMK0901-00202}{Mahasiswa mengerjakan UTS yang mencakup konsep teknologi, inovasi teknologi, IPTEK, etika rekayasa, perkembangan ICT, dan sistem komputer.}{Mengerjakan soal UTS secara mandiri melalui platform ujian online dalam waktu 60 menit.}{Jawaban UTS online yang terekam pada platform ujian.}{Penguasaan konsep TI, inovasi, IPTEK, etika, dan ICT & 7 & Menjawab benar penerapan konsep TI, inovasi, IPTEK, etika, dan perbedaan ICT/TIK pada kasus yang diberikan. & Konsep dasar dijawab benar, tetapi penerapannya pada kasus belum lengkap. & Jawaban tidak menunjukkan pemahaman konsep TI dan ICT. \\ \\
Penjelasan struktur dan komponen sistem komputer & 4 & Menjelaskan struktur komputer, perangkat I/O, register, memory, CPU, CU, ALU, dan BUS beserta fungsinya dengan benar. & Menyebutkan komponen komputer dengan benar, tetapi sebagian fungsi tertukar. & Tidak mampu menyebutkan komponen sistem komputer beserta fungsinya. \\ \\
Penjelasan cara kerja dan arsitektur sistem komputer & 4 & Menjelaskan elemen, arsitektur, dan interkoneksi antar komponen sistem komputer secara runtut. & Menjelaskan arsitektur secara umum, tetapi hubungan antar komponen belum runtut. & Tidak mampu menjelaskan cara kerja dan arsitektur sistem komputer. \\}

\assessmentblock{Kuis 2 Komputasi dan Logika}{Kuis (Ujian Online)}{SCPMK0901-00203}{Mahasiswa mengerjakan kuis online tentang representasi data, aljabar Boolean, dan flowchart.}{Mengerjakan soal kuis secara mandiri melalui platform ujian online dalam waktu 60 menit.}{Jawaban kuis online yang terekam pada platform ujian.}{Ketepatan konversi sistem bilangan dan representasi data & 5 & Mengonversi bilangan antar basis dan menentukan representasi/tipe data dengan langkah dan hasil benar. & Langkah konversi benar, tetapi terdapat kesalahan hitung pada sebagian soal. & Tidak memahami langkah konversi bilangan atau representasi data. \\ \\
Ketepatan operasi logika dan penyederhanaan ekspresi Boolean & 4 & Menyelesaikan operasi logika, tabel kebenaran, dan penyederhanaan ekspresi Boolean sesuai hukum aljabar Boolean. & Operasi logika dasar benar, tetapi penyederhanaan ekspresi belum tuntas. & Operasi logika dan ekspresi Boolean dijawab keliru. \\ \\
Ketepatan penyusunan flowchart untuk kasus sederhana & 3 & Menyusun flowchart dengan simbol yang benar dan alur logika yang menyelesaikan kasus. & Simbol flowchart benar, tetapi alur logika belum sepenuhnya menyelesaikan kasus. & Simbol dan alur flowchart tidak sesuai kaidah atau tidak menjawab kasus. \\}

\assessmentblock{UAS Konsep Teknologi Informasi}{UAS (Ujian Online)}{SCPMK0901-00203, SCPMK0208-00204, SCPMK0502-00205}{Mahasiswa mengerjakan UAS yang mencakup representasi data, aljabar Boolean, flowchart, jaringan komputer dan internet, aplikasi TI, sertifikasi bidang TI, dan aspek keamanan dasar sistem operasi/jaringan/cloud.}{Mengerjakan soal UAS secara mandiri melalui platform ujian online dalam waktu 60 menit.}{Jawaban UAS online yang terekam pada platform ujian.}{Penerapan representasi data, aljabar Boolean, dan flowchart pada kasus & 15 & Menyelesaikan kasus komputasi dengan konversi bilangan, ekspresi Boolean, dan flowchart yang benar dan runtut. & Sebagian penerapan benar, tetapi terdapat kesalahan langkah atau hasil pada beberapa kasus. & Penerapan komputasi keliru atau tidak menunjukkan proses penyelesaian. \\ \\
Pemahaman jaringan komputer dan internet & 10 & Menjelaskan jenis jaringan, topologi, perangkat jaringan, serta perbedaan internet dan intranet sesuai kebutuhan kasus. & Menyebutkan jenis dan perangkat jaringan dengan benar, tetapi pemilihannya pada kasus kurang tepat. & Tidak mampu membedakan jenis jaringan, topologi, dan perangkat jaringan. \\ \\
Pemilihan aplikasi TI dan sertifikasi sesuai kebutuhan bidang & 7 & Memilih aplikasi TI dan sertifikasi bidang TI yang sesuai kebutuhan pengguna/organisasi disertai alasan yang tepat. & Pilihan aplikasi atau sertifikasi relevan, tetapi alasan pemilihannya masih umum. & Pilihan aplikasi TI dan sertifikasi tidak relevan dengan kebutuhan kasus. \\ \\
Pengenalan aspek keamanan dasar sistem operasi, jaringan, dan cloud & 3 & Mengidentifikasi ancaman dan mekanisme keamanan dasar (firewall, autentikasi) pada sistem operasi, jaringan, dan platform cloud dengan tepat. & Mengenali sebagian ancaman/mekanisme keamanan dasar, tetapi penjelasannya masih umum. & Tidak mampu mengidentifikasi aspek keamanan dasar yang relevan. \\}

\begin{tblr}{width=\textwidth,colspec={|Q[l,m,3.2cm]|X[l,m]|},hlines,vlines,hspan=minimal,row{1-3}={bg=headgray},rowsep=1.5pt,colsep=3pt}
Jadwal Pelaksanaan: & Minggu 1--17 sesuai RPS. Setiap evaluasi memiliki RTM dan rubrik multi-kriteria. \\
Lain-Lain Yang Diperlukan: & LMS, platform ujian online, perangkat lunak sesuai mata kuliah, dan perangkat presentasi. \\
Pustaka: & Mengacu pustaka utama dan pendukung pada bagian RPS. \\
\end{tblr}
